\section{An uncertainty barrier}\label{yo:sec:barrier}

Every prime side computation is limited in one of two ways. Compact support in the level variable keeps the prime sum finite but limits the level to $\log n\le X$. A finite basis limits the frequency, the ordinate $|\tau|\le T$. For the tent coefficient $D$ of Proposition~\ref{yo:prop:tent} the limit in level runs into zero verification, which is what the next proposition makes precise.

\begin{proposition}\label{yo:prop:barrier}
Suppose all zeros with $0<\gamma\le H$ lie on the critical line, and let $D_{\rm on}\ge0$ be the part of $D$ from on line zeros. Then for every $x>0$
\[
|D(x)-D_{\rm on}(x)|\le\varepsilon_H(x):=8\cosh^2\frac x4\sum_{\gamma>H}\frac1{\gamma^2},
\]
the sum being over all zeros above $H$, counted with multiplicity. With $H=3\times10^{12}$ (Platt and Trudgian \cite{PT21}, the height $H_0$ of Theorem~\ref{thm:dbeta:minall}) the sum differs from $(\log(H/2\pi)+1)/(2\pi H)=1.4797\times10^{-12}$ by less than $2\times10^{-24}$ (Theorem~\ref{ag:thm:explicit}), so that $\varepsilon_H(x)<1.184\times10^{-11}\cosh^2(x/4)$.
\end{proposition}

\begin{proof}
Off line zeros have $|\gamma|>H$ and $|\delta|<\frac12$, so $|\sin(\gamma_\rho x/2)|^2\le\cosh^2(\delta x/2)\le\cosh^2(x/4)$ and $|\gamma_\rho|^2\ge\gamma^2$. Each off line zero therefore contributes at most $4\cosh^2(x/4)/\gamma^2$ to $D$, and summing over all such zeros, two (a zero and its conjugate) for each ordinate $\gamma>H$ of the upper half plane, gives the bound.
\end{proof}

The proposition is of use only against a lower bound for $D_{\rm on}$. Every term of $D_{\rm on}$ is nonnegative, so the first $N$ zeros, which are on the line, give one: $D_{\rm on}(x)\ge L_N(x):=\sum_{n\le N}8\sin^2(\gamma_nx/2)/\gamma_n^2$. For the $N=1700$ certified zeros of the book the least value of $L_N$ on $[0.05,47.1]$ is $0.040704$, at $x=25.383$, and $0.04070$ is a lower bound for it computed by branch and bound (\texttt{yo\_tent.py}, \texttt{code/\allowbreak README.md}); on $[0.05,16]$ the least value is $0.0531$.

\begin{proposition}[Computed positivity of the tent coefficient]\label{yo:prop:Dpos}
For the $N=1700$ certified zeros of the book let $L_N(x)=\sum_{n\le N}8\sin^2(\gamma_nx/2)/\gamma_n^2$, and let $H=3\times10^{12}$. Since all zeros with $0<\gamma\le H$ lie on the critical line \cite{PT21}, $D(x)\ge L_N(x)-\varepsilon_H(x)$ for every $x>0$. Moreover $L_N(x)>\varepsilon_H(x)$ on $[0.05,47.48]$, so $D>0$ on $[0.05,47.48]$ from the verification of RH up to $H$ and the first $1700$ zeros alone, with no prime sum. The first inequality is proved. The second is computed in double precision and not in ball arithmetic: for $x<46.688$ it follows from the floor $0.04070$ of $L_N$ on $[0.05,47.1]$, found by branch and bound, and for $46.688\le x\le47.48$ from a grid of step $5\times10^{-4}$ and the Lipschitz constant $10.85$ of $L_N$.
\end{proposition}

The level $47.10$ at which $\varepsilon_H$ reaches $0.05$ ($e^{47.10}=2.85\times10^{20}$) is arithmetically right, but it is the barrier only if $D_{\rm on}$ stays above $0.05$, which it does not; with the floor $0.0407$ the bound reaches it at $x=46.69$ ($e^{46.69}=1.9\times10^{20}$), and pointwise at $x=47.48$ ($e^{47.48}=4.2\times10^{20}$). The barrier is at about $47$ either way, because $\varepsilon_H$ grows exponentially. A prime side computation of $D$ at a level below $47.48$, which needs all prime powers up to $e^x$, can therefore only confirm a positive sign that the zeros already give; it cannot produce a negative value of $D$, which is how an off line zero would show. To see one through $D$ a computation must reach levels beyond $47.48$, with all prime powers up to $4\times10^{20}$, or heights beyond $H$, and the second is the zero verification itself.

In frequency, Proposition~\ref{yo:prop:mult} makes the same point for a basepoint $\alpha\ge\sigma_{\max}$: a test vector with spectrum concentrated in $|\tau|\le T<H$ sees only the symbol there, and a zero $\beta+i\gamma$ with $\gamma>H$ contributes to the Hadamard sum for the symbol the term $(\alpha-\beta)/((\alpha-\beta)^2+(\tau-\gamma)^2)$, of modulus at most $|\alpha-\beta|/(\gamma-T)^2$.

What the barrier does not say. Proposition~\ref{yo:prop:barrier} is proved for the single coefficient $D$ and for no other vector. A normalised function supported in the window $[-a,a]$ and modulated to a height above $H$ (Proposition~\ref{yo:prop:W}) sees an off line zero there with a weight that does not decay with the height, and $\varepsilon_H$, which carries the factor $\gamma^{-2}$ of the tent, does not bound the contribution of the zeros above $H$ to the form of such a vector. Nor could a perturbation bound of size $10^{-11}$ settle the sign of a windowed form at $\alpha=\frac12$, whose least eigenvalue is below $10^{-15}$ from $a=0.8$ on (Table~\ref{yo:tab:lam}). The tent is one vector of the window, so the level $x_*$ of Remark~\ref{yo:rem:scale} bounds Yoshida's threshold only from above, $a_0\le x_*/2$, and the window experiment of Chapter~\ref{chap:zhu} (Table~\ref{zh:tab:det}) finds a hypothetical off line pair at half lengths far below $x_*/2$.

\section{The question of a path to RH}\label{yo:sec:path}

The chapter does not give a path to RH. Every criterion proved here has the same structure. On the RH side it is a positivity statement: $D\ge0$, $T_k\ge0$ for infinitely many $k$, the modulated heat trace of Proposition~\ref{yo:prop:W} nonnegative, a Pick function, a coercive damped form. On the other side it is the presence of a pole off the real axis, or a negative Cauchy kernel, at an off line ordinate. The barrier of Section~\ref{yo:sec:barrier} then says, for the tent coefficient $D$, that a computation at levels below $47$ cannot go beyond what zero verification already knows; for the other quantities no such statement is proved there. A proof needs an inequality valid at all levels or all heights at once, and it must use the Euler product: the function $f_{\rm DH}$ of Davenport and Heilbronn \cite{DH36} (Section~\ref{pk:sec:theta}) shares the functional equation and has zeros off the line, so it violates the analogue of every criterion in this chapter.

What the chapter does change is the shape of the missing inequality. Theorem~\ref{yo:thm:coercive} puts it in its sharpest form:
\begin{equation}\label{yo:eq:target}
\re\frac{\xi'}{\xi}(\alpha+i\tau)\ \ge\ \frac{\xi'}{\xi}(\alpha)\qquad\text{for all }\tau,\ \alpha\in(\tfrac12,1],
\end{equation}
which we write $v_\alpha(\tau)\ge b_\alpha$ as in Theorem~\ref{thm:dbeta:minall}. Taken as a family over $\alpha\in(\frac12,1]$, and not line by line, this is equivalent to the coercivity of every damped Yoshida form with the constant it has on functions of low frequency. Three facts about it:
\begin{enumerate}
\item It is true with no hypothesis for every $\alpha\ge1$ and every $\tau$ (Propositions~\ref{bp:prop:v1min} and \ref{bp:prop:valpha}), and for $\frac12<\alpha\le\alpha_{\max}=7.35$ and $0<|\tau|\le H-\frac{1+\sqrt2}2$ (Theorem~\ref{thm:dbeta:minall}(ii)); for $\frac12<\alpha<1$ and all $\tau$ it is equivalent to RH (Theorem~\ref{thm:dbeta:minall}(vi)). On the line $\alpha=1$ the minimum is at $\tau=0$, the next lowest local minimum of $v_1$, at $\tau=17.49$, is $4.90$ times $b_1$, and on $[20,2000]$ the least value is $8.23\,b_1$, at $\tau=27.67$ (computed by \texttt{yo\_symbol.py}, \texttt{code/\allowbreak README.md}).
\item If it holds, even in the weak form $v_{\alpha'}(\tau)>0$, on every line $\re s=\alpha'$ with $\alpha<\alpha'\le1$, then $\sigma_{\max}\le\alpha$, since a zero $\rho_0=\beta_0+i\gamma_0$ with $\beta_0>\alpha$ sends $v_{\alpha'}(\gamma_0)$ to $-\infty$ as $\alpha'\uparrow\beta_0$; on all of $(\frac12,1]$ it is RH. On a single line it is not known to exclude a zero to its right (Remark~\ref{rem:ch2:oneline}). Its $\alpha\downarrow\frac12$ limit, under RH, is the sharpened Laguerre inequality $-(\log\Xi)''(\tau)\ge2p_1$.
\item Its failure mode is local. A zero with $\beta>\alpha$ would create a negative Cauchy kernel of height $1/(\pi(\beta-\alpha))$ at its ordinate, competing against a symbol of size $\frac1{2\pi}\log\tau$ there. So the inequality at $\alpha$ asks for control of $-\re\zeta'/\zeta(\alpha+i\tau)$ by $\frac12\log(\tau/2\pi)$, a one sided bound on $\zeta'/\zeta$; the Fej\'er tapered form of the question is Proposition~\ref{mw:prop:AX}. No method known today proves it at any fixed $\alpha<1$ for all $\tau$; the zero free half plane $\re s>\frac78$ announced in October 2026 (Section~\ref{sec:ch25:open}) has not been refereed and is not counted.
\end{enumerate}

\section{Open problems}\label{yo:sec:open}

What remains:

\begin{problem}[The symbol inequality below the line $\alpha=1$]\label{yo:op:symbol}
The inequality $v_\alpha(\tau)\ge b_\alpha$ is a theorem of the book for every $\alpha\ge1$ and every $\tau$ (Propositions~\ref{bp:prop:v1min} and \ref{bp:prop:valpha}), and for $\frac12<\alpha\le\alpha_{\max}$ up to height $3\times10^{12}-1.21$ (Theorem~\ref{thm:dbeta:minall}(ii)); for $\frac12<\alpha<1$ and all $\tau$ it is equivalent to RH (Theorem~\ref{thm:dbeta:minall}(vi)). What is open is the region $\alpha<1$ above the verified height. Determine the part of it in which $v_\alpha(\tau)>b_\alpha$ can be proved with no hypothesis and with explicit constants; the Vinogradov Korobov region $\alpha\ge1-c(\log|\tau|)^{-2/3}(\log\log|\tau|)^{-1/3}$ \cite{Vinogradov58,Korobov58} is the natural target. A proof on every line of a strip $\alpha_0<\alpha\le1$ and for all $\tau$ would give $\sigma_{\max}\le\alpha_0$ (Section~\ref{yo:sec:path}), which is why the target is a region that narrows with the height.
\end{problem}

\begin{problem}[A prime side proof at $\alpha=1$]\label{yo:op:primes}
The book proves $v_1(\tau)\ge b_1$ through the zeros (Proposition~\ref{bp:prop:v1min}), by the Cauchy kernels of the ordinates nearest to $\tau$ and a bound on the gaps between them (compare Lemma~\ref{yo:lem:gap}). At $\alpha=1$ the statement can be written entirely in primes. Since $\xi'/\xi(1)=1+\frac{\gamma_E}2-\frac12\log4\pi$ and $\re\frac1{i\tau}=0$, it reads
\[
\sum_{n\ge2}\frac{\Lambda(n)}{n}\cos(\tau\log n)\ \le\ \frac12\re\psi\Bigl(\frac{1+i\tau}2\Bigr)-\frac12\log\pi+\frac1{1+\tau^2}-\frac{\xi'}{\xi}(1)\qquad(\tau\ne0),
\]
the left side being the conditionally convergent series for $-\re\zeta'/\zeta(1+i\tau)$. Find a proof using only Mertens type information, in the spirit of the Fej\'er taper of Proposition~\ref{mw:prop:AX}; the right side grows like $\frac12\log(\tau/2\pi)$ and the left side is $O(\log\log\tau)$ under RH.
\end{problem}

\begin{problem}[L\"owner at a single window]\label{yo:op:lowner}
By Theorem~\ref{yo:thm:lowner}, positivity on $K(a)$ is positivity of the L\"owner matrices of $g_a$ on the lattice $\frac\pi a\mathbb Z$. L\"owner's hard theorem \cite{Lowner34,Donoghue74} recovers the Pick property from L\"owner matrices on all finite subsets of an interval. What does positivity on the lattice alone, for one fixed $a$, imply about the location of the zeros? A positive answer that grew with $a$ would quantify Yoshida's threshold $a_0$.
\end{problem}

\begin{problem}[The constants of the ladder and of the window]\label{yo:op:orders}
Determine the order of $K_0$ for the ladder and of $a_0$ for Yoshida's window, for a single off line pair, and compare with the displacement $\sigma_{\max}-\frac12$ and with the de Bruijn Newman constant $\Lambda\le2(\sigma_{\max}-\frac12)^2$ (de Bruijn \cite{deBruijn}, in the normalisation of Rodgers and Tao \cite{RodgersTao} used in Section~\ref{ch:dbn}). For the ladder the computations of Section~\ref{yo:sec:rungs} give $K_0(\sec\theta_*-1)\approx0.70$ for the pair of the stress test (Remark~\ref{yo:rem:stress}) and, over fourteen pairs, a first failing order that grows like $\gamma^2$ and, at fixed $\gamma$, like $1/\delta$ and not like the heuristic $\gamma^2/\delta^2$; its dependence on $\delta$ and on the spacing of the true ordinates is what is to be determined. For the window only $a_0\le x_*/2$ is known (Remark~\ref{yo:rem:scale}), and the window experiment of Chapter~\ref{chap:zhu} (Table~\ref{zh:tab:det}) finds the threshold far below it.
\end{problem}

\begin{problem}[Zero density and the barrier]\label{yo:op:density}
Proposition~\ref{yo:prop:barrier} allows every zero above $H$ to have $\delta$ arbitrarily close to $\frac12$. Zero density estimates make zeros with large $\delta$ rare. Sharpen the barrier level using an explicit $N(\sigma,T)$ bound, and determine whether any prime side computation of $D$ at feasible levels can, even in principle, exclude off line zeros beyond the verified height. For windows no barrier is proved: decide whether a window of half length $a$ below about $23.7$ (level $47.48$) can see an off line zero above $H$, for which Section~\ref{zh:sec:horizon} of Chapter~\ref{chap:zhu} gives numerical evidence and a heuristic horizon, not a proof.
\end{problem}

\subsection*{Notes}

The forms of this chapter are Yoshida's \cite{Yoshida92}, and the explicit formula is Weil's \cite{Weil52}. Lagarias's criterion \cite{Lagarias} asks only for $v_\alpha>0$; the minimum at the origin of Theorem~\ref{thm:dbeta:minall}(vi) sharpens it. In integrated form, under RH, the minimum at the origin says that $\sigma\mapsto|\xi(\sigma+i\tau)|/\xi(\sigma)$ is nondecreasing on $(\frac12,\infty)$, a normalised version of the reformulation of Sondow and Dumitrescu \cite{SD10} that $|\xi(\sigma+i\tau)|$ increases in $\sigma$. Theorem~\ref{yo:thm:D} is classical in substance.
